% Part: intuitionistic-logic
% Chapter: tableaux
% Section: rules

\documentclass[../../../include/open-logic-section]{subfiles}

\begin{document}

\olfileid{int}{tab}{rul}

\olsection{د شهودي منطق قاعدې}

د $\land$ او $\lor$ نښلوونکو قاعدې د عادي قضييزو نښه‌لرونکو
!!{tableau}s له قاعدو سره يو شان دي؛ يوازې مخکښونه ورزيات شوي
دي۔ په هر حالت کښې پر نښه‌لرونکي !!{formula}
$\sFmla{S}{!A}[\sigma]$ پلې شوې قاعده داسې نوي
!!{formula}s جوړوي چې مخکښ يې هم~$\sigma$ وي۔ دا له مفهومي
پلوه څرګنده ده: مثلاً که $!A \land !B$ په~$\sigma$ نومولې
نړۍ کښې رښتيا وي، نو $!A$ او $!B$ هم په~$\sigma$ کښې رښتيا
دي (په بلې نړۍ کښې اړين نۀ دي)۔ د $\land$ او $\lor$ قاعدې
په \olref{tab:prop-rules} کښې راټولوو۔

\begin{table}
  \[\def\arraystretch{3}\begin{array}{|c|c|}
    \hline
    \AxiomC{\sFmla{\True}{!A \land !B}[\sigma]}
    \RightLabel{\TRule{\True}{\land}}
    \UnaryInfC{\sFmla{\True}{!A}[\sigma]}
    \noLine
    \UnaryInfC{\sFmla{\True}{!B}[\sigma]}
    \DisplayProof
    &
    \AxiomC{\sFmla{\False}{!A \land !B}[\sigma]}
    \RightLabel{\TRule{\False}{\land}}
    \UnaryInfC{$\sFmla{\False}{!A}[\sigma] \quad \mid \quad
      \sFmla{\False}{!B}[\sigma]$}
    \DisplayProof
    \\[2ex]
    \hline
    \AxiomC{\sFmla{\True}{!A \lor !B}[\sigma]}
    \RightLabel{\TRule{\True}{\lor}}
    \UnaryInfC{$\sFmla{\True}{!A}[\sigma] \quad \mid \quad
      \sFmla{\True}{!B}[\sigma]$}
    \DisplayProof
    &
    \AxiomC{\sFmla{\False}{!A \lor !B}[\sigma]}
    \RightLabel{\TRule{\False}{\lor}}
    \UnaryInfC{\sFmla{\False}{!A}[\sigma]}
    \noLine
    \UnaryInfC{\sFmla{\False}{!B}[\sigma]}
    \DisplayProof
    \\[2ex]
    \hline
  \end{array}\]
  \caption{د $\land$ او $\lor$ لپاره د مخکښ‌لرونکي !!{tableau} قاعدې}
  \ollabel{tab:prop-rules}
\end{table}

د تړل کېدو شرط د عادي !!{tableau}s له شرط سره ورته دے، خو
يوازې د !!{formula}s برابري کافي نۀ ده؛ مخکښونه هم بايد سره
برابر يا د لاسرسي له مخې اړوند وي۔ په حقيقت کښې لږ پراخ شرط
اخيستلے شو: په شهودي مدلونو کښې !!{formula}s چې يو ځل رښتيا
شي، بيا رښتيا پاتې کېږي؛ نو ناشونې ده چې $!A$ په~$\sigma$
کښې رښتيا خو په کوم لاس‌رسي وړ مخکښ~$\sigma.{*}$ کښې کاذب
وي۔ نو څانګه هله تړلې وي چې دواړه لاندې نښې ولري:
\[
\sFmla{\True}{!A}[\sigma] \quad\text{او}\quad \sFmla{\False}{!A}[\sigma.{*}]
\]
د کوم مخکښ~$\sigma$ او !!{formula}~$!A$ دپاره۔ پام وکړئ چې
که نښې سرچپه وي، يعنې څانګه دا ولري:
\[
\sFmla{\False}{!A}[\sigma] \quad\text{او}\quad \sFmla{\True}{!A}[\sigma.{*}]
\]
نو څانګه يوازې هغه وخت تړلې ده چې $*$ تشه لړۍ وي۔

پر دې سربېره، که څانګه~$\sFmla{\True}{\bot}[\sigma]$ ولري،
هغه هم تړلې ده۔

د فرضونو د ايښودلو قاعدې هم د عادي !!{tableau}s په شان دي،
خو د فرضونو دپاره تل مخکښ~$1$ کاروو۔ (کوم مخکښ چې کاروو مهم
نۀ دے، بس د ټولو فرضونو دپاره يو شان وي۔) نو مثلاً وايو چې
\[
!B_1, \dots, !B_n \Proves !A
\]
هله او يوازې هله چې د لاندې فرضونو دپاره تړلے تابلو وي:
\[
\sFmla{\True}{!B_1}[1], \dots, \sFmla{\True}{!B_n}[1],
\sFmla{\False}{!A}[1].
\]

د شرط~$\lif$ قاعدې له کلاسيکو او مودالي قاعدو بېلې دي۔ د
$\TRule{\True}{\lif}$ قاعده هغه څانګه چې
$\sFmla{\True}{!A \lif !B}[\sigma]$ لري، په دوو بېلو
څانګو کښې په ترتيب $\sFmla{\False}{!A}[\sigma.{*}]$ يا
$\sFmla{\True}{!B}[\sigma.{*}]$ ورزياتولو غځوي۔ د ژباړې د
نسخې څرګندونه (OLINT-013): منجمده سرچينه دلته د ريښتيني
شرط د قاعدې پايلې د کاذب شرط د شاهد په بڼه سرچپه ليکي؛ د
همدې برخې لاندې جدول د ريښتيني شرط دپاره کاذب مقدم او رښتينې
پايله په دوو څانګو کښې ښيي، او همدا دلته مراد دي۔ قاعده يوازې
پر داسې مخکښ~$\sigma.{*}$ پلې کېدے شي چې په اړونده څانګه کښې
\emph{لا مخکې} راغلے وي۔ داسې مخکښ «کارېدلے» بولو۔ (څرنګه چې
$\sigma.{*}$ خپله $\sigma$ هم رانغاړي، قاعده تل د
$\sFmla{\False}{!A}[\sigma]$ او
$\sFmla{\True}{!B}[\sigma]$ په بېلو څانګو کښې په زياتولو
پلې کېدے شي۔) د نسخې څرګندونه (OLINT-023): په منجمده سرچينه کښې د شرط د دواړو قاعدو په نثري نومونو کښې د نښې او نښلوونکي د ماکرو آرګومېنټونه سرچپه دي؛ دلته د همدې برخې د جدول او د ماکرو د تعريف له مخې سم شوي دي۔

د $\TRule{\False}{\lif}$ قاعده هغه څانګه چې
$\sFmla{\False}{!A \lif !B}[\sigma]$ لري، په هماغه څانګه
کښې د $\sFmla{\True}{!A}[\sigma.n]$ او
$\sFmla{\False}{!B}[\sigma.n]$ دواړو په زياتولو غځوي؛
$\sigma.n$ بايد په څانګه کښې نوے مخکښ وي۔

د $\lnot$ قاعدې په ورته ډول تعريفېږي (د $\lnot !A$ د
$!A \lif \lfalse$ په توګه تعريف کاروو)۔

قاعدې په \olref{tab:rules-lif-lnot} کښې ورکړل شوې دي۔

\begin{table}
  \[\def\arraystretch{3}\begin{array}{|c|c|}
    \hline
    \AxiomC{\sFmla{\True}{\lnot !A}[\sigma]}
    \RightLabel{\TRule{\True}{\lnot}}
    \UnaryInfC{$\sFmla{\False}{!A}[\sigma.{*}]$}
    \DisplayProof
    &
    \AxiomC{\sFmla{\False}{\lnot !A}[\sigma]}
    \RightLabel{\TRule{\False}{\lnot}}
    \UnaryInfC{\sFmla{\True}{!A}[\sigma.n]}
    \DisplayProof
    \\[1ex]
    \text{$\sigma.{*}$ کارېدلے دے} & \text{$\sigma.n$ نوے دے}\\
    \hline
    \AxiomC{\sFmla{\True}{!A \lif !B}[\sigma]}
    \RightLabel{\TRule{\True}{\lif}}
    \UnaryInfC{$\sFmla{\False}{!A}[\sigma.{*}] \quad \mid \quad
      \sFmla{\True}{!B}[\sigma.{*}]$}
    \DisplayProof
    &
    \AxiomC{\sFmla{\False}{!A \lif !B}[\sigma]}
    \RightLabel{\TRule{\False}{\lif}}
    \UnaryInfC{\sFmla{\True}{!A}[\sigma.n]}
    \noLine
    \UnaryInfC{\sFmla{\False}{!B}[\sigma.n]}
    \DisplayProof
    \\[1ex]
    \text{$\sigma.{*}$ کارېدلے دے} & \text{$\sigma.n$ نوے دے}\\
    \hline
  \end{array}\]
  \caption{د $\lnot$ او $\lif$ لپاره د مخکښ‌لرونکي !!{tableau} قاعدې}
  \ollabel{tab:rules-lif-lnot}
\end{table}

\end{document}
